\documentclass[conference]{IEEEtran}
\IEEEoverridecommandlockouts
% The preceding line is only needed to identify funding in the first footnote. If that is unneeded, please comment it out.
\usepackage{cite}
\usepackage{amsmath,amssymb,amsfonts}
\usepackage{algorithmic}
\usepackage{graphicx}
\usepackage{textcomp}
\usepackage{xcolor}
\usepackage{amsmath}
\usepackage{amssymb}
\usepackage{bbm}
\usepackage{ebproof}
\usepackage{babel}
\usepackage{graphicx}
\usepackage{hyperref}
\usepackage{tikz}
\usepackage{backnaur}
\usepackage{tikz-cd}
\usetikzlibrary{cd}
\usepackage{agda}
\usepackage{import}

\def\BibTeX{{\rm B\kern-.05em{\sc i\kern-.025em b}\kern-.08em
    T\kern-.1667em\lower.7ex\hbox{E}\kern-.125emX}}
\begin{document}

\title{Schémas composables non biaisés pour les combinateurs catégoriques et combinateurs catégoriques cartésiens.}

\author{\IEEEauthorblockN{Eliott Houget}
\IEEEauthorblockA{\textit{Équipe Cosynus} \\
\textit{LIX, CNRS, École polytechnique, Institut Polytechnique de
	Paris}\\
Palaiseau, France \\
eliott.houget@ens-rennes.fr}
}
\maketitle

\begin{abstract}Nous proposons une adaptation de la notion de schéma composable pour les catégories non biaisées et les catégories cartésiennes. Cela est une première étape pour proposer une variante à la construction de la théorie des types $CCCaTT_1$ telle qu'introduite par S. Mimram. Nous présentons les combinateurs catégoriques ainsi qu'une relation d'équivalence $\sim$ qui les identifie relativement aux axiomes de la théories des catégories. Nous introduisons un nouveau système de réécriture pour définir une forme normalisée des termes. Nous exposons une définition inductive des schémas composables et montrons que ceux-ci sont contractibles à équivalence $\sim$ près. Enfin, nous exposons comment adapter les notions précédentes après avoir introduit un produit et un objet terminal. Une partie de ce travail est implémentée en Agda.
\end{abstract}

\begin{IEEEkeywords}
Théorie des types, théorie des catégories, réécriture
\end{IEEEkeywords}

\section{Introduction}

\subsection{Théorie des types} La théorie des types est une branche de la logique qui s'attache à formaliser des termes mathématiques typés. Formellement, une théorie des types est un langage formel défini à l'aide d'un ensemble de règles de réécriture. Du point de vue de la logique, les types apparaissent comme des propositions (et les types dépendants comme des prédicats), tandis que les termes sont des preuves de leur type. Les preuves (i.e les termes) d'une théorie des types sont par définition toujours constructives (et construites comme des arbres dont les nœuds sont les règles du système de réécriture). Ainsi, elles peuvent être vérifiées mécaniquement à l'aide d'outils d'assistance de preuve tels qu'Agda, Rocq (anciennement Coq), Lean, etc. Les définitions, propositions et preuves de ce rapports sont formalisées et implémentés en Agda ; ainsi la théorie des types est ici autant un outil qu'un objet d'étude.

\subsection{Théorie des catégories} La théorie des catégories est un vaste domaine décrivant les structures abstraites des mathématiques. Les catégories sont composées d'objets et de flèches entre ces derniers ; l'ajouts d'une notion de produit d'objets et d'un objet terminal définit les catégories cartésiennes.

\subsection{Motivation} L'objectif à long terme est la construction d'une théorie des types dépendants non biaisée pour les catégories cartésiennes closes de dimension supérieure, dans laquelle la notion d'égalité est remplacée par des témoins d'égalité (de dimension supérieure). Ce but s'inscrit lui-même dans un projet d'ériger des fondations constructivistes des mathématiques. Pour cela S. Mimram introduit la théorie des types non-baisée $CCCaTT_1$. Elle est dite non-biaisée car il existe une unique règle d'introduction des termes. Cette règle s'appuie sur un calcul de séquent minimal introduisant les schémas composables. Ils définissent un ensemble de types contractiles (types uniquement habités). Cet ensemble est de taille assez grande pour permettre de dérivés un certain nombre de séquents essentiels pour montrer la cohérence de $CCCaTT_1$ vis-à-vis du $\lambda$-calcul simplement typé. De plus, la définition des schémas composables est plus aisement manipulable que la définition des types contractiles.

\subsection{Plan du papier} Ce rapport se découpe en trois parties. Dans la partie II, nous présenterons brièvement le cadre, les notions et les outils nécessaires à la bonne compréhension de ce travail. Les parties III et IV présentent chacune une construction et une preuve de la bonne définition des schémas composables dans le cas des catégories, d'une part, et dans le cas des catégories cartésiennes, d'autre part.

\section{Cadres, notions, notations et outils}

\subsection{Catégories} Une catégorie est la donnée d'une collection $\mathcal{C}_0$ d'objets, d'une collection $\mathcal{C}_1$ de morphismes (ou flèches) telles que : 
\begin{itemize}
	\item pour chaque morphisme $f$ il existe un objet $s(f)$ (appelé source) et un objet $t(f)$ (appelé but);
	\item pour chaque paire de morphismes $f$ et $g$ tels que $t(f) = s(g)$ il existe un morphisme $g \cdot f$ (appelé composé de $f$ et $g$);
	\item pour chaque objet $X$, il existe un morphisme $id_X$ (appelé identité de $X$);
	\item les quatre propriétés suivantes sont vérifiées :
	\begin{itemize}
		\item la source et le but sont compatibles avec la composition : $s(g \cdot f) = s(f)$ et $t(g \cdot f) = t(g)$;
		\item la source et le but sont compatibles avec les identités : $s(id_X) = X$ et $t(id_X) = X$;
		\item Si $t(f) = s(g)$ et $t(g) = s(h)$, alors $(h \cdot g) \cdot f = h \cdot (g \cdot f)$;
		\item Si $s(f) = X$ et $t(f) = Y$, alors $id_X \cdot f = f = f \cdot id_Y$.
	\end{itemize}
\end{itemize}

\subsection{Catégories cartésiennes} Une catégorie $\mathcal{C}$ est dite cartésienne, si elle est munie d'un objet terminal $\mathbbm 1$ et si pour toute paire d'objets $A$ et $B$, il existe un produit $A \times B$ dans $\mathcal C$.
Dans une catégorie cartésienne $\mathcal{C}$, pour toute paire d'objets $A$ et $B$, il existe deux morphismes de projection : $$\pi_0 : A \times B \rightarrow A$$ et $$\pi_1 : A \times B \rightarrow A$$. De plus, pour toute paire de morphisme $f$ et $g$ telle que $s(f) = s(g)$, il existe un unique morphisme $\left<f, g\right>$ tel que le diagramme suivant commute :

\begin{center}
	\begin{tikzcd}
		& X \arrow[d, "\left<f\, g\right>"] \arrow[rd, "g"] \arrow[ld, "f"]\\
		A
		& A \times B \arrow[r, "\pi_1"] \arrow[l, "\pi_0"]
		& B
	\end{tikzcd}
\end{center}

Ainsi les quatre propriétés suivantes sont vérifiées :
\begin{itemize}
	\item $\pi_0 \circ \left<f, g\right> = f$;
	\item $\pi_1 \circ \left<f, g\right> = g$;
	\item $\left<\pi_0, \pi_1\right> = id$;
	\item $\left<g , h\right> \circ f = \left<g \circ f, h \circ f\right>$
\end{itemize}

\subsection{Combinateurs catégoriques typés} Une théorie des types catégoriques est une théorie des types dont les termes sont les morphismes d'une catégorie $\mathcal{C}$ typés par leur source et leur cible (appartenant à $\mathcal{C}_0$).

\subsection{Schémas composables} Ces schémas composables prennent la forme d'ensembles de flèches entre des types, ces flèches sont telles qu'elles sont composables d'une manière unique. 

\section{Pasting scheme dans des catégories simples}

On suppose fixée une petite catégorie $\mathcal{C}$.

\subsection{Types} L'ensemble des objets $\mathcal{C}_0  = \left\{x, y, \dots\right\}$ définit un ensemble de variables de type. Une \textbf{flèche} est un couple de variables de types, notée $x \rightarrow y$, avec $x$ la \textbf{source} et $y$ la \textbf{cible}. Ainsi tout morphisme $f$ de $\mathcal{C}_1$ est typé en notant $f : s(f) \rightarrow t(f)$.

\subsection{Contextes} Les contextes prennent la forme de listes de flèches définies inductivement par : $$ \Gamma ::= \epsilon \ \text{(contexte vide)} \ |\ \Gamma , f : x \rightarrow y$$
On notera $f : A \in \Gamma$ pour signifier que $f : A$ est une flèche apparaissant au moins une fois dans le contexte $\Gamma$. On note $Con$ l'ensemble des contextes. De plus, on note $t(\Gamma)$ l'ensemble des cibles des flèches de $\Gamma$.


\subsection{Termes} Un terme typé est un morphisme $f$ typé par une flèche $A$ et construit dans un contexte $\Gamma$. On définit inductivement le séquent $\Gamma \vdash f : x \rightarrow y$ par les règles suivantes :
\begin{center}
	\begin{prooftree}
		\hypo[]{f : A \in \Gamma}
		\infer1[(var)]{\Gamma \vdash f : A}
	\end{prooftree}
	\hspace{10pt}
	\begin{prooftree}
		\infer0[(id)]{\Gamma \vdash id_x : x \rightarrow x}
	\end{prooftree} \\
	\vspace{20pt}
	\begin{prooftree}
		\hypo[]{\Gamma \vdash f : x \rightarrow y}
		\hypo[]{\Gamma \vdash g : y \rightarrow z}
		\infer2[(comp)]{\Gamma \vdash f \cdot g : x \rightarrow z}
	\end{prooftree}
\end{center}
On note $Tm(\Gamma, A)$ l'ensemble des termes de contexte $\Gamma$ et de flèche $A$.

\subsection{Relation $\sim$}
On définit la relation $\sim\ \subset Tm(\Gamma, A)$ qui identifie les termes égaux relativement aux quatre axiomes de la définition d'une catégorie présentée dans la partie II. 
\begin{center}
	\begin{prooftree}
		\infer0[(unitl)]{f \cdot id \sim f}
	\end{prooftree}
	\hspace{10pt}
	\begin{prooftree}
		\infer0[(unitr)]{id \cdot f \sim f}
	\end{prooftree} \\
	\vspace{20pt}
	\begin{prooftree}
		\infer0[(assoc)]{f \cdot (g \cdot h) \sim (f \cdot g) \cdot h}
	\end{prooftree} \\
	\vspace{20pt}
	\begin{prooftree}
		\hypo[]{f_l \sim f_r}
		\hypo[]{g_l \sim g_r}
		\infer2[(comp)]{f_l \cdot g_l \sim f_r \cdot g_r}
	\end{prooftree} \\
	\vspace{20pt}
	\begin{prooftree}
		\infer0[($\sim$-refl)]{f \sim f}
	\end{prooftree}
	\hspace{20pt}
	\begin{prooftree}
		\hypo[]{f \sim g}
		\infer1[($\sim$-sym)]{g \sim f}
	\end{prooftree} \\
	\vspace{20pt}
	\begin{prooftree}
		\hypo[]{f \sim g}
		\hypo[]{g \sim h}
		\infer2[($\sim$-trans)]{f \sim h}
	\end{prooftree}
\end{center}

Par les règles $\sim$-refl, $\sim$-sym et $\sim$-trans, $\sim$ est une relation d'équivalence.

\subsection{Termes normaux}

On souhaite normaliser les arbres de dérivation des termes en les transformant en des peignes. Ces termes normalisés ont pour but d'être des représentants canoniques pour les classes d'équivalence de la relation $\sim$. 
On introduit un systeme d'écriture composé de 2 règles engendrant les termes normalisés. Ce système de réécriture doit apparaître comme un intermédiaire entre le système très souple qui définit les termes, et le système de réécriture linéaire qui définira les schémas composables.
\begin{center}
	\begin{prooftree}
		\infer0[norm-id]{\Gamma \vdash_{Norm} id_x : x \rightarrow x}
	\end{prooftree} \\
	\vspace{20pt}
	\begin{prooftree}
		\hypo[]{\Gamma \vdash_{Norm} f : x \rightarrow y}
		\hypo[]{y \rightarrow z \in \Gamma}
		\infer2[norm-comp]{\Gamma \vdash_{Norm} f \cdot var_{y \rightarrow z} : x \rightarrow z}
	\end{prooftree} \\
\end{center}

\paragraph*{Composition de deux termes normaux} Afin de définir la transformation de normalisation, on introduit une transformation auxiliaire. C'est par cette transformation que la mise sous forme de peigne des arbres de dérivation devient effective.

\begin{equation}
\begin{gathered}
	M \left(
	\begin{prooftree}
		\hypo[]{\Pi_f}
		\infer1[]{\Gamma \vdash_{Norm} f : x \rightarrow y}
	\end{prooftree}
	,
	\begin{prooftree}
		\infer0[]{\Gamma \vdash_{Norm} id_y : y \rightarrow y}
	\end{prooftree}
	\right) \\
	:= \\
	\begin{prooftree}
		\hypo[]{\Pi_f}
		\infer1[]{\Gamma \vdash_{Norm} f : x \rightarrow y}
	\end{prooftree}
\end{gathered}
\end{equation}

\begin{equation}
\begin{gathered}
	M \left(
	\begin{prooftree}
		\hypo[]{\Pi_f}
		\infer1[]{\Gamma \vdash_{Norm} f : w \rightarrow x}
	\end{prooftree}
	,
	\begin{prooftree}
		\hypo[]{\Pi_g}
		\infer1[]{\Gamma \vdash_{Norm} g : x \rightarrow y}
		\hypo[]{y \rightarrow z \in \Gamma}
		\infer2[]{\Gamma \vdash_{Norm} g \cdot var_{y \rightarrow z} : x \rightarrow z}
	\end{prooftree}
	\right) \\
	:= \\
	\begin{prooftree}
		\hypo[]{
			M \left(
			\begin{prooftree}
				\hypo[]{\Pi_f}
				\infer1[]{\Gamma \vdash_{Norm} f : w \rightarrow x}
			\end{prooftree}
			,
			\begin{prooftree}
				\hypo[]{\Pi_g}
				\infer1[]{\Gamma \vdash_{Norm} g : x \rightarrow y}
			\end{prooftree}
			\right)
		}
		\hypo[]{y \rightarrow z \in \Gamma}
		\infer2[]{\Gamma \vdash_{Norm} (f \cdot g) \cdot var_{y \rightarrow z} : x \rightarrow z}
	\end{prooftree}
\end{gathered}
\end{equation}


\paragraph*{Normalisation} On définit une transformation $N$ de la dérivation d'un terme renvoyant une dérivation du terme normal correspondant:
\begin{equation}
	N \left(
	\begin{prooftree}
		\infer0[]{\Gamma \vdash id_x : x \rightarrow x}
	\end{prooftree}
	\right)
	:=
	\begin{prooftree}
		\infer0[]{\Gamma \vdash_{Norm} id_x : x \rightarrow x}
	\end{prooftree}
\end{equation}

\begin{equation}
\begin{gathered}
	N \left(
	\begin{prooftree}
		\hypo[]{x \rightarrow y \in \Gamma}
		\infer1[]{\Gamma \vdash var_{x \rightarrow y} : x \rightarrow y}
	\end{prooftree}
	\right) \\
	:= \\
	\begin{prooftree}
		\infer0[]{\Gamma \vdash_{Norm} id_x : x \rightarrow x}
		\hypo[]{x \rightarrow y \in \Gamma}
		\infer2[]{\Gamma \vdash_{Norm} id_x \cdot var_{x \rightarrow y} : x \rightarrow y}
	\end{prooftree}\\
\end{gathered}
\end{equation}

\begin{equation}
\begin{gathered}
	N \left(
	\begin{prooftree}
		\hypo[]{\Pi_f}
		\infer1[]{\Gamma \vdash f : x \rightarrow y}
		\hypo[]{\Pi_g}
		\infer1[]{\Gamma \vdash g : y \rightarrow z}
		\infer2[]{\Gamma \vdash f \cdot g : y \rightarrow z}
	\end{prooftree}
	\right) \\
	:= \\
	M \left(
	N \left(
	\begin{prooftree}
		\hypo[]{\Pi_f}
		\infer1[]{\Gamma \vdash f : x \rightarrow y}
	\end{prooftree}
	\right),
	N \left(
	\begin{prooftree}
		\hypo[]{\Pi_g}
		\infer1[]{\Gamma \vdash g : y \rightarrow z}
	\end{prooftree}
	\right)
	\right)
\end{gathered}
\end{equation}

\textbf{lemme} Soient $f$ et $g$ dans $Tm(\Gamma, A)$, si $N(f) \equiv N(g)$, alors $f \sim g$.

\subsection{Schémas composables} Un schéma composable est un couple contexte/flèche $(\Gamma, A)$ tels qu'il existe exactement un unique terme $\Gamma \vdash f : A$ à $\sim$-équivalence près.

On introduit le calcul de séquent suivant :
\begin{center}
	\begin{prooftree}
		\infer0[(ax)]{\emptyset \vdash_{ps} x \rightarrow x}
	\end{prooftree}
	\hspace{50pt}
	\begin{prooftree}
		\hypo[]{\Gamma \vdash_{ps} x \rightarrow y}
		\infer1[(ext) $z \notin b(\Gamma) \cup \{x\}$]{\Gamma , y \rightarrow z \vdash_{ps} x \rightarrow z}
	\end{prooftree}
\end{center} 

Nous montrons que ses deux règles définissent $PS \subset Con \times Arr$ un ensemble de schéma composable. Formellement, $\forall \Gamma \in Con, \forall (x , y) \in Ty^2, ((\Gamma, x \rightarrow y) \in PS \Rightarrow (\forall (f, g) \in Tm(\Gamma, x \rightarrow y)^2, f \sim g))$.

\textbf{lemme} Soient $\Gamma$ un contexte, $A$ une flèche, $\Gamma \vdash_{Norm} f : A$ et $\Gamma \vdash_{Norm} g : A$ deux termes normaux. Si $(\Gamma, A)$ est un schéma composable alors $f \equiv g$. \\

\textbf{Theoreme} Soient $\Gamma$ un contexte, $A$ une flèche, $\Gamma \vdash f : A$ et $\Gamma \vdash g : A$ deux termes. Si $(\Gamma, A)$ est un schéma composable alors $f \sim g$.

\textbf{\textit{preuve}} On utilise les deux lemmes précédents : $N(f)$ et $N(g)$ sont deux termes normaux dont le couple contexte/flèche est un schéma composable ainsi, par le lemme 2, $N(f) \equiv N(g)$. Subséquemment, par le lemme 1, $f \sim g$.

\subsection{Existence}
Par le théorème précédent nous avons obtenu l'unicité sous réserve d'existence. Nous montrons que chaque schéma composable est habité par un terme. Pour cela nous construisons inductivement un terme en déconstruisant la l'arbre de dérivation du schéma composable. Formellement, pour tout context $\Gamma$, toute flèche $A$ tel que $\Gamma \vdash_{ps} A$ est dérivable, il existe un morphisme $f$ tel que $\Gamma \vdash f : A$ est dérivable.


\subsection{Séquents}
On peut donc dériver les séquents suivants, indispensables dans la construction de $CCCaTT_1$.

\begin{enumerate}
	\item $\emptyset \vdash_{ps} x \rightarrow x$
	\item $x \rightarrow y \vdash_{ps} x \rightarrow y$
	\item $f : x \rightarrow y, g : y \rightarrow z \vdash_{ps} x \rightarrow z$
	\item $f : x \rightarrow y \vdash id \cdot f = f$
	\item $f : x \rightarrow y \vdash f \cdot id = f$
	\item $f : x \rightarrow y, g : y \rightarrow z, h : z \rightarrow w \vdash (f \cdot g) \cdot h = f \cdot (g \cdot h)$
\end{enumerate}

\section{Pasting scheme dans des catégories cartésiennes}
Nous n'avons pas pu conclure la preuve formelle implémentée en Agda, nous présentons ici les grandes idées de la preuve telle qu'envisagée. Nous détaillons les objets introduits ainsi que la défintion de schéma composable proposée. Nous expliquerons ce qui manque pour conclure.
Nous enrichissons la définition des types en ajoutant les deux règles suivantes :
\begin{center}
	\begin{prooftree}
		\infer0[$({\mathbbm 1}_I)$]{\mathbbm 1 \in Ty}
	\end{prooftree}
	\hspace{20pt}
	\begin{prooftree}
		\hypo[]{A \in Ty}
		\hypo[]{B \in Ty}
		\infer2[$(\times_I)$]{A \times B \in Ty}
	\end{prooftree}
\end{center}

Afin d'exprimer qu'une variable de type apparait dans un type produit nous introduisons le prédicat $\blacktriangleright$ défini inductivement de la manière suivante :
\begin{center}
	\begin{prooftree}
		\infer0[$(\blacktriangleright_{ax})$]{x \blacktriangleright x}
	\end{prooftree}
	\hspace{30pt}
	\begin{prooftree}
		\hypo[]{A \blacktriangleright x}
		\infer1[$(\blacktriangleright_g)$]{A \times B \blacktriangleright x}
	\end{prooftree}
	\hspace{30pt}
	\begin{prooftree}
		\hypo[]{B \blacktriangleright x}
		\infer1[$(\blacktriangleright_d)$]{A \times B \blacktriangleright x}
	\end{prooftree}
\end{center}

\paragraph{Termes} Pour $\Gamma$ un contexte et $A$ une flèche nous enrichissons la définition de $Tm(\Gamma, A)$ :
\begin{center}
	\begin{prooftree}
		\infer0[(term)]{\Gamma \vdash term_x : x \rightarrow \mathbbm 1}
	\end{prooftree} \\
	\vspace{20pt}
	\begin{prooftree}
		\hypo[]{\Gamma \vdash f : x \rightarrow y}
		\hypo[]{\Gamma \vdash g : x \rightarrow z}
		\infer2[(pair)]{\Gamma \vdash \left<f, g\right> : x \rightarrow y \times z}
	\end{prooftree} \\
	\vspace{20pt}
	\begin{prooftree}
		\infer0[(fst)]{\Gamma \vdash \pi_0 : x \times y \rightarrow x}
	\end{prooftree}
	\hspace{20pt}
	\begin{prooftree}
		\infer0[(snd)]{\Gamma \vdash \pi_1 : x \times y \rightarrow y}
	\end{prooftree}
\end{center}

\subsection{Contexte} Sans perte de généralité, nous travaillons sous l'hypothèse que toute flèche appartenant à un contexte a une cible qui n'est ni un produit, ni un objet terminal. Cette hypothèse n'est pas contraignante car un morphisme produisant une paire peut être vu comme une pair de morphisme (dont la source est identique).

\subsection{Relation $\sim$}
De même, nous enrichissons la relation d'équivalence $\sim$ afin d'identifier les termes selon les axiomes des catégories cartésiennes.

\begin{center}
	\begin{prooftree}
		\hypo[]{\Gamma \vdash f : x \rightarrow y}
		\hypo[]{\Gamma \vdash g : x \rightarrow z}
		\infer2[(pfst)]{\left<f, g\right> \cdot \pi_0 \sim f}
	\end{prooftree} \\
	\vspace{20pt}
	\begin{prooftree}
		\hypo[]{\Gamma \vdash f : x \rightarrow y}
		\hypo[]{\Gamma \vdash g : x \rightarrow z}
		\infer2[(psnd)]{\left<f, g\right> \cdot \pi_1 \sim g}
	\end{prooftree} \\
	\vspace{20pt}
	\begin{prooftree}
		\infer0[(pext)]{\left<\pi_0, \pi_1\right> \sim id}
	\end{prooftree} \\
	\vspace{20pt}
	\begin{prooftree}
		\hypo[]{\Gamma \vdash f : x \rightarrow y}
		\hypo[]{\Gamma \vdash g : y \rightarrow z}
		\hypo[]{\Gamma \vdash h : y \rightarrow z}
		\infer3[(pnat)]{f \cdot \left<g, g\right> \sim \left<f \cdot g, f \cdot h\right>}
	\end{prooftree}
\end{center}

\subsection{Termes normaux}
Comme nous l'avons fait dans la partie II, on introduit un système d'écriture engendrant des termes normalisés. Ces termes normalisés ont pour but d'être des représentants canoniques pour les classes d'équivalence de la relation $\sim$.

\begin{center}
	\begin{prooftree}
		\infer0[norm-term]{\Gamma \vdash_{Norm} term_A : A \rightarrow \mathbbm 1}
	\end{prooftree} \\
	\vspace{20pt}
	\begin{prooftree}
		\hypo[]{A \blacktriangleright x}
		\infer1[norm-proj]{\Gamma \vdash_{Norm} proj_{A,x} : A \rightarrow x}
	\end{prooftree} \\
	\vspace{20pt}
	\begin{prooftree}
		\hypo[]{\Gamma \vdash_{Norm} t : A \rightarrow B}
		\hypo[]{B \rightarrow x \in \Gamma}
		\infer2[norm-var]{\Gamma \vdash_{Norm} t \cdot var_{B \rightarrow x} : A \rightarrow x}
	\end{prooftree} \\
	\vspace{20pt}
	\begin{prooftree}
		\hypo[]{\Gamma \vdash_{Norm} f : A \rightarrow B}
		\hypo[]{\Gamma \vdash_{Norm} g : A \rightarrow C}
		\infer2[norm-pair]{\Gamma \vdash_{Norm} \left<f, g\right> : A \rightarrow B \times C}
	\end{prooftree}
\end{center}

\paragraph{Affaiblissements de la source d'un terme normal} Avec C un élément de $Ty$, on définit les transformations $\mathrm{Aff}_g(\cdot, C)$ et $\mathrm{Aff}_d(\cdot, C)$ qui étant donné une preuve de $\Gamma \vdash_{Norm} f : A \rightarrow B$ renvoit respectivement $\Gamma \vdash_{Norm} f : A \times C \rightarrow B$ et $\Gamma \vdash_{Norm} f : C \times A \rightarrow B$.

\begin{align}
\begin{gathered}
	\mathrm{Aff}_g \left(
	\begin{prooftree}
		\infer0[]{\Gamma \vdash_{Norm} term_A : A \rightarrow \mathbbm 1}
	\end{prooftree}
	, C
	\right) \\
	:= \\
	\begin{prooftree}
		\infer0[]{\Gamma \vdash_{Norm} term_{A \times C} : A \times C \rightarrow \mathbbm 1}
	\end{prooftree}
\end{gathered}
\end{align}

\begin{align}
\begin{gathered}
	\mathrm{Aff}_g \left(
	\begin{prooftree}
		\hypo[]{A \blacktriangleright x}
		\infer1[]{\Gamma \vdash_{Norm} proj_{A,x} : A \rightarrow x}
	\end{prooftree} 
	, C
	\right) \\
	:= \\
	\begin{prooftree}
		\hypo[]{A \blacktriangleright x}
		\infer1[$\blacktriangleright_g$]{A \times C \blacktriangleright x}
		\infer1[]{\Gamma \vdash_{Norm} proj_{A \times C,x} : A \times C \rightarrow x}
	\end{prooftree}
\end{gathered}
\end{align}

\begin{align}
\begin{gathered}
	\mathrm{Aff}_g \left(
	\begin{prooftree}
		\hypo[]{\Gamma \vdash_{Norm} f : A \rightarrow B}
		\hypo[]{B \rightarrow x \in \Gamma}
		\infer2[]{\Gamma \vdash_{Norm} f \cdot var_{B \rightarrow x} : A \rightarrow x}
	\end{prooftree}
	, C
	\right) \\
	:= \\
	\begin{prooftree}
		\hypo[]{\mathrm{Aff}_g \left(\Gamma \vdash_{Norm} f : A \rightarrow B , C\right)}
		\hypo[]{B \rightarrow x \in \Gamma}
		\infer2[]{\Gamma \vdash_{Norm} f \cdot var_{B \rightarrow x} : A \times C \rightarrow x}
	\end{prooftree}
\end{gathered}
\end{align}

\begin{align}
\begin{gathered}
	\mathrm{Aff}_g \left(
	\begin{prooftree}
		\hypo[]{\Gamma \vdash_{Norm} f : A \rightarrow B}
		\hypo[]{\Gamma \vdash_{Norm} g : A \rightarrow C}
		\infer2[]{\Gamma \vdash_{Norm} \left<f, g\right> : A \rightarrow B \times C}
	\end{prooftree}
	, D
	\right) \\
	:= \\
	\begin{prooftree}
		\hypo[]{\mathrm{Aff}_g \left( \Gamma \vdash_{Norm} f : A \rightarrow B , D \right)}
		\hypo[]{\mathrm{Aff}_g \left( \Gamma \vdash_{Norm} g : A \rightarrow C , D \right)}
		\infer2[]{\Gamma \vdash_{Norm} \left<f, g\right> : A \times D \rightarrow B \times C}
	\end{prooftree}
\end{gathered}
\end{align}

$\mathrm{Aff}_d$ est définie de manière similaire.

\paragraph{Normalisation} On construit une fonction de normalisation N qui à un terme associe un terme normal en modifant inductivement la preuve : 
\begin{align}
	N\left( 
	\begin{prooftree}
		\hypo[]{\Gamma \vdash_{Norm} f : A \rightarrow B}
		\hypo[]{\Gamma \vdash_{Norm} g : A \rightarrow C}
		\infer2[norm-pair]{\Gamma \vdash_{Norm} \left<f, g\right> : A \rightarrow B \times C}
	\end{prooftree}
	\right)
\end{align}

\begin{enumerate}
	\item $\emptyset \vdash_{ps} A \rightarrow \mathbbm 1$
	\item $f : A \rightarrow B, g : A \rightarrow C \vdash_{ps} A \rightarrow B \times C$
	\item $\emptyset \vdash_{ps} A \times B \rightarrow A$
	\item $\emptyset \vdash_{ps} A \times B \rightarrow B$
	
	\item $f : A \rightarrow B, g : A \rightarrow C \vdash \left<f, g\right> \cdot fst = f$
	\item $f : A \rightarrow B, g : A \rightarrow C \vdash \left<f, g\right> \cdot snd = g$
	\item $f : A \rightarrow B, g : B \rightarrow C, h : C \rightarrow D \vdash f \cdot \left<g, h\right> = \left< f \cdot g, f \cdot h\right>$
	\item $\emptyset \vdash \left<fst, snd\right> = id$
	\item $f : A \rightarrow B \vdash f \cdot term_B = term_A$
\end{enumerate}

\begin{thebibliography}{00}
\bibitem{b1} Samuel Mimram, ``An unbiased simply typed combinatory logic,'' LIX, CNRS, École polytechnique, Institut Polytechnique de Paris, Palaiseau, France, 2026
\bibitem{b2} Dominique Boucher, ``La Théorie des catégories en informatique :
notions de base et application,'' Département d'informatique et recherche opérationnelle
Faculté des études supérieures, Université de Montréal, Août 1993
\bibitem{b3} Alexander Grothendieck, Section 4 of: Techniques de construction et théorèmes d’existence en géométrie algébrique III: préschémas quotients, Séminaire Bourbaki: années 1960/61, exposés 205-222, Séminaire Bourbaki, no. 6 (1961), Exposé no. 212
\end{thebibliography}

\appendices

\section{Preuves séquents catégories simples}

\begin{enumerate}
	\item
	\begin{prooftree}
		\infer0[(ax)]{\emptyset \vdash_{ps} x \rightarrow x}
	\end{prooftree}
	\vspace{15pt}
	\item
	\begin{prooftree}
		\infer0[(ax)]{\emptyset \vdash_{ps} x \rightarrow x}
		\infer1[(ext)]{x \rightarrow y \vdash_{ps} x \rightarrow y}
	\end{prooftree}
	\vspace{15pt}
	\item
	\begin{prooftree}
		\infer0[(ax)]{\emptyset \vdash_{ps} x \rightarrow x}
		\infer1[(ext)]{x \rightarrow y \vdash_{ps} x \rightarrow y}
		\infer1[(ext)]{x \rightarrow y, y \rightarrow z \vdash_{ps} x \rightarrow z}
	\end{prooftree}
	\vspace{15pt}
	\item
	\begin{prooftree}
		\infer0[(ax)]{\emptyset \vdash_{ps} x \rightarrow x}
		\infer1[(ext)]{x \rightarrow y \vdash_{ps} x \rightarrow y}
		\hypo[]{\Pi_{id \cdot f}}
		\hypo[]{f : x \rightarrow y \in f : x \rightarrow y}
			\infer1[(var)]{f : x \rightarrow y \vdash f : x \rightarrow y}
		\infer3[$(ps^=)$]{f : x \rightarrow y \vdash id \cdot f = f}
	\end{prooftree}
	avec
	\begin{align*}
	\begin{gathered}
	\Pi_{id \cdot f} \\
	:= \\
	\begin{prooftree}
		\infer0[id]{f : x \rightarrow y \vdash id_x : x \rightarrow x}
		\hypo[]{f : x \rightarrow y \in f : x \rightarrow y}
		\infer1[var]{f : x \rightarrow y \vdash f : x \rightarrow y}
		\infer2[comp]{f : x \rightarrow y \vdash id_x \cdot f : x \rightarrow y}
	\end{prooftree}
	\end{gathered}
	\end{align*}
	
	\item
	\begin{prooftree}
		\infer0[ax]{\emptyset \vdash_{ps} x \rightarrow x}
		\infer1[ext]{x \rightarrow y \vdash_{ps} x \rightarrow y}
		\hypo[]{\Pi_{f \cdot id}}
		\hypo[]{f : x \rightarrow y \in f : x \rightarrow y}
		\infer1[var]{f : x \rightarrow y \vdash f : x \rightarrow y}
		\infer3[$(ps^=)$]{f : x \rightarrow y \vdash f \cdot id = f}
	\end{prooftree}
	avec
	\begin{align*}
		\begin{gathered}
			\Pi_{f \cdot id} \\
			:= \\
			\begin{prooftree}
				\hypo[]{f : x \rightarrow y \in f : x \rightarrow y}
				\infer1[var]{f : x \rightarrow y \vdash f : x \rightarrow y}
				\infer0[id]{f : x \rightarrow y \vdash id_x : x \rightarrow x}
				\infer2[comp]{f : x \rightarrow y \vdash id_x \cdot f : x \rightarrow y}
			\end{prooftree}
		\end{gathered}
	\end{align*}
	\vspace{15pt}
	\item
	\begin{prooftree}
		\hypo[]{\Pi_{ps}}
		\hypo[]{\Pi_{(f \cdot g) \cdot h}}
		\hypo[]{\Pi_{f \cdot (g \cdot h)}}
		\infer3[$(ps^=)$]{\Gamma \vdash (f \cdot g) \cdot h = f \cdot (g \cdot h)}
	\end{prooftree} \\
	\vspace{15pt}
	avec \\
	$\Gamma = f : x \rightarrow y, g : y \rightarrow z, h : z \rightarrow w$ \\ 
	\begin{align*}
		\begin{gathered}
			\Pi_{id \cdot f} \\
			:= \\
			\begin{prooftree}
				\infer0[ax]{\emptyset \vdash_{ps} x \rightarrow x}
				\infer1[ext]{x \rightarrow y \vdash_{ps} x \rightarrow y}
				\infer1[(ext)]{x \rightarrow y, y \rightarrow z \vdash_{ps} x \rightarrow z}
				\infer1[(ext)]{x \rightarrow y, y \rightarrow z, z \rightarrow w \vdash_{ps} x \rightarrow w}
			\end{prooftree}
		\end{gathered}
	\end{align*} \\
	\begin{align*}
		\begin{gathered}
			\Pi_{id \cdot f} \\
			:= \\
			\begin{prooftree}
				\hypo[]{f : x \rightarrow y \in \Gamma}
				\infer1[(var)]{\Gamma \vdash f : x \rightarrow y}
				\hypo[]{g : y \rightarrow z \in \Gamma}
				\infer1[(var)]{\Gamma \vdash g : y \rightarrow z}
				\hypo[]{h : z \rightarrow w \in \Gamma}
				\infer1[(var)]{\Gamma \vdash h : z \rightarrow w}
				\infer2[(comp)]{\Gamma \vdash g \cdot h : y \rightarrow w}
				\infer2[(comp)]{\Gamma \vdash (f \cdot g) \cdot h  : x \rightarrow w}
			\end{prooftree}
		\end{gathered}
	\end{align*} \\
	\begin{align*}
		\begin{gathered}
			\Pi_{(f \cdot g) \cdot h} \\
			:= \\
			\begin{prooftree}
				\hypo[]{f : x \rightarrow y \in \Gamma}
				\infer1[(var)]{\Gamma \vdash f : x \rightarrow y}
				\hypo[]{g : y \rightarrow z \in \Gamma}
				\infer1[(var)]{\Gamma \vdash g : y \rightarrow z}
				\infer2[(comp)]{\Gamma \vdash f \cdot g : x \rightarrow z}
				\hypo[]{h : z \rightarrow w \in \Gamma}
				\infer1[(var)]{\Gamma \vdash h : z \rightarrow w}
				\infer2[(comp)]{\Gamma \vdash (f \cdot g) \cdot h  : x \rightarrow w}
			\end{prooftree}
		\end{gathered}
	\end{align*}
\end{enumerate}

\section{Preuves séquents catégories cartésiennes}


\end{document}
